\section{Introduction}
\label{sec:introduction}

Stokes' theorem is one of the places where the informal compression power of
mathematics is most visible.  In a textbook proof on a compactly supported
form, one takes a finite chart cover, chooses a partition of unity, applies the
Euclidean theorem on each coordinate box, observes that the artificial faces
cancel or vanish, and identifies the remaining terms with the boundary
integral \cite{spivak1965calculus,munkres1991analysis,lee2013smooth,tu2011introduction}.
The entire local-to-global passage may occupy a page.  In a proof
assistant, however, none of these choices is silent.  The finite cover has an
index type.  The support of each partition coefficient has to be controlled.
The chart representatives are total functions.  Boundary charts are modeled on
a half-space rather than on all of \(\R^{n+1}\).  The auxiliary faces of
coordinate boxes must be separated from the true boundary face.  The final
finite sum must be reconstructed as a canonical global endpoint.

This paper presents a Lean 4 formalization of that local-to-global core for a
coordinate-represented compact-support Stokes theorem, building on Lean's
dependent type theory and mathlib's analysis, topology, and manifold
infrastructure \cite{demoura2015lean,demoura2021lean4,leanprover,tpil4,mathlib2020,mathlibdocs}.
The main public theorem is the first-principles theorem named:
\begin{lstlisting}[language=Lean]
CoverIndexedZeroCompactRepresentedStokesFirstPrinciplesInput.representedStokes
\end{lstlisting}
Its public input consists of chartwise smoothness of the \(n\)-form and
compactness of the closure of its support.  From these two analytic hypotheses,
under the ambient half-space model assumptions carried by the surrounding Lean
context, the development constructs the hidden data of the proof: pointwise
chart boxes, an open finite selected cover, a support-controlled partition, a
boundary half-space box refinement, smooth refined coefficients,
zero-localized chart representatives, local half-space Stokes fields, and two
canonical represented endpoints.  It then proves that the generated bulk
endpoint equals the generated boundary endpoint.

The theorem is represented rather than native in a precise sense.  The
endpoints are not arbitrary real numbers supplied by a user.  They are the
finite coordinate sums produced by the classical proof after localizing the
form and applying Stokes on finitely many half-space boxes.  The theorem proves
the equality
\[
  \BulkRep(\omega) = \BoundaryRep(\omega),
\]
where both sides are generated by the formal construction.  What is not claimed
is a comparison theorem between these represented endpoints and a future
mathlib-native API for integrating differential forms over oriented manifolds
with boundary.  That comparison is important, but it is a different layer: it
identifies a canonical finite coordinate construction with a chosen native
global integral.  The present artifact proves the Stokes proof payload at the
represented level.

\subsection{What Lean forces us to expose}

The traditional statement
\[
  \int_M \dd\omega = \int_{\partial M} \omega
\]
suppresses several constructions.  The formalization shows that these are not
mere implementation details.

First, compact support is a constructive resource.  The proof does not simply
use compactness abstractly; it uses compactness to extract a finite family of
chart boxes and to make every local summand subordinate to one of them.  The
development therefore works with a compact carrier
\[
  K = \overline{\supp(\omega)}
\]
and proves that the selected partition has topological support controlled over
\(K\).  This support control is later used to kill artificial boundary faces.

Second, boundary charts are half-space charts.  A coordinate box adjacent to
the boundary is not an ordinary closed box sitting in an ambient open subset of
\(\R^{n+1}\).  It touches the face \(x_0=0\), and the local theorem must
separate the open-box differentiability needed for the bulk term from the
closed-box continuity and face regularity needed for boundary terms.  This is
why the development contains a genuine half-space local Stokes layer, centered
around \code{HalfSpaceBoxInteriorStokesFields}, rather than pretending that
boundary boxes have ordinary Euclidean neighborhoods across the boundary.

Third, Lean functions are total.  A chart representative such as
\code{transitionPullbackInChart} has values even outside the meaningful chart
source, target, or overlap.  A naive theorem saying that the support of this
total representative is contained in a selected coordinate box is often the
wrong theorem.  The project introduces
\code{ManifoldForm.transitionPullbackInChartZero}, a zero-localized chart
representative, and proves support theorems for that representative.  The key
formal statement
\[
\code{zeroSupportOfSelectedSmoothRefinement}
\]
turns the informal phrase ``the localized form is supported in the box'' into a
Lean-checkable topological support inclusion for the actual refined
representatives used by the global proof.

Fourth, the endpoint reconstruction is itself mathematics.  Once each local
half-space box satisfies Stokes and the artificial faces vanish, one still has
to prove that the nested finite sum over selected boundary charts and refined
pieces is the canonical represented bulk endpoint, and similarly for the
boundary endpoint.  This is not a naming convention.  It is a finite-sum
reconstruction theorem built into the localized refined endpoint layer.

\subsection{Scale of the artifact}

The active project contains 544 Lean files under \code{Stokes/}, totaling about
\(1.65\times 10^5\) lines.  The size is not accidental.  A large part of the
work is the removal of theorem-facing certificates.  Earlier intermediate
routes exposed selected covers, refinement data, local Stokes records,
endpoint adapters, target-image control, or support assumptions.  These
intermediate APIs were useful for stabilizing the proof, but the final public
route compresses them behind first-principles input.  The default public import
has also been slimmed so that the preferred entry point is the first-principles
represented theorem rather than the historical development scaffolding.

The theorem is audited in the usual Lean way.  The repository contains no
project-level uses of \code{sorry}, \code{admit}, or new axioms in the active
Lean sources, and the theorem-level axiom audit for the first-principles
statement reports the standard Lean/mathlib axioms
\[
\code{propext},\quad \code{Classical.choice},\quad \code{Quot.sound}.
\]
The paper records the commands used for this audit in
\cref{sec:artifact-scale,app:theorem-map}.

\subsection{Certificate generation as a contribution}

A useful way to read the development is as a transformation from
certificate-consuming formalization to certificate-generating formalization.
An early formal milestone for a theorem of this kind can be logically correct
while still asking the user for proof objects that arise only inside the
local-to-global argument: a selected finite cover, collar boxes, refined
partition coefficients, local Stokes fields, support-in-box assumptions,
image-control packages, and endpoint adapters.  Such a theorem proves a real
identity, but its interface is not yet the mathematical theorem one wants to
state.

The final route continues past that milestone.  Each former certificate field
is either proved unnecessary or generated from earlier mathematical data.  The
zero-localized representative is the sharpest example.  The tempting support
assumption for the ordinary total chart representative is not merely
inconvenient; it is the wrong shape of theorem.  The proof replaces the
representative by one that agrees on the chart source and is zero elsewhere,
then proves the topological support statement required by half-space Stokes.
This is a mathematical correction of the formal object, not a cosmetic API
change.

\subsection{Contributions}

The main contributions are the following.

\begin{itemize}
  \item A coordinate-represented compact-support Stokes theorem in Lean 4 whose
  public input is chartwise smoothness and compactness of the closed support.

  \item A formal local-to-global construction turning compact support into an
  open finite selected cover, a support-controlled partition, and a smooth
  boundary half-space refinement.

  \item A half-space local Stokes layer that separates open-interior
  differentiability, closed-box regularity, artificial-face vanishing, and the
  outward-first sign of the true boundary face.

  \item A zero-localization method for chart representatives, replacing false
  ordinary-support assumptions for totalized chart functions by generated
  topological-support theorems for the correct local representative.

  \item A canonical represented endpoint layer proving equality of generated
  finite coordinate bulk and boundary sums, rather than consuming those
  endpoints as external certificates.

  \item A proof-engineering compression from certificate-consuming
  intermediate routes to a certificate-generating first-principles public
  theorem.
\end{itemize}

\subsection{Organization}

\Cref{sec:represented-stokes} explains represented endpoints and the exact
claim of the theorem.  \Cref{sec:proof-decomposition} gives the local-to-global
proof decomposition from the perspective of formalization.  The next four
sections discuss the main mathematical bottlenecks: half-space local Stokes
(\cref{sec:halfspace}), zero-localized support (\cref{sec:zero-localization}),
compact finite data (\cref{sec:compact-finite-data}), and endpoint
reconstruction (\cref{sec:endpoint-reconstruction}).  \Cref{sec:main-theorem}
states the main Lean theorem and describes the first-principles route.
\Cref{sec:artifact-scale} records artifact scale and verification.  Related
work and limitations appear in \cref{sec:related-work,sec:limitations}.  The
appendix gives a theorem and API map for readers navigating the Lean files.
